\clearpage
\thispagestyle{plain}
\section*{Einleitung des Autors}

\begin{center}
{\arabicfont\fontsize{15}{20}\selectfont\begin{Arabic}بِسْمِ اللَّهِ الرَّحْمَنِ الرَّحِيمِ\end{Arabic}}
\end{center}

\begin{center}
{\arabicfont\fontsize{13}{18}\selectfont\begin{Arabic}
الحمد لله رب العالمين، والعاقبة للمتقين، وصلى الله وسلم على عبده ورسوله نبينا محمد، وعلى آله
وأصحابه أجمعين.

أما بعد: فهذه كلمات موجزة في بيان بعض ما يجب أن يعرفه العامة عن دين الإسلام، وسميتها:
﴾الدروس المهمة لعامة الأمة﴿.

وأسأل الله أن ينفع بها المسلمين، وأن يتقبلها مني، إنه جواد كريم.
\end{Arabic}}
\end{center}

\begin{flushright}
{\arabicfont\begin{Arabic}عبد العزيز بن عبد الله بن باز\end{Arabic}}
\end{flushright}

\vspace{1.5em}
\begin{center}\color{MutedText}\rule{28mm}{0.4pt}\end{center}
\vspace{1em}

Im Namen Allahs, des Allerbarmers, des Barmherzigen.

Alles Lob gebührt Allah, dem Herrn der Weltenbewohner, und das gute Ende gehört den
Gottesfürchtigen, und Allahs Segen und Frieden seien auf Seinem Diener und Gesandten, unserem
Propheten Muhammad, und auf seiner Familie und all seinen Gefährten.

Um fortzufahren: Im Folgenden sind einige zusammengefasste Worte, die darlegen, was die
Allgemeinheit der Menschen über die Religion des Islams wissen sollte. Ich gab dieser Schrift den
Namen „Die wichtigen Lektionen für die allgemeine muslimische Gemeinschaft" (\emph{ad-Durūs
al-Muhimmah li-ʿĀmmati l-Ummah}).

Und ich bitte Allah, dass Er den Muslimen Nutzen daraus gewährt und dass Er es von mir annimmt.
Wahrlich, Er ist freigebig und edelmütig.

\vspace{1em}
\begin{flushright}
\textbf{ʿAbd al-ʿAzīz ibn ʿAbdillāh ibn Bāz} \ar{رحمه الله}
\end{flushright}

\vspace{0.5em}
{\small\color{MutedText}Arabischer Urtext nach \texttt{sources/durus-al-muhimmah-arabisch-urtext-scan.pdf},
Seite 3. Deutsche Übersetzung nach \texttt{addurus-almuhimmah-ar-de.pdf}, S.\,4.}
\clearpage
